\documentclass[11pt,a4paper]{article}
\usepackage[utf8]{inputenc} % Not needed with XeLaTeX
\usepackage{fontspec}
\usepackage[margin=0.75in]{geometry}
\usepackage{xcolor}
\definecolor{darknavy}{RGB}{0,0,100}
\usepackage[colorlinks=true, urlcolor=darknavy, linkcolor=darknavy]{hyperref}
\usepackage{enumitem}
% \usepackage{fontawesome}
\usepackage{fontawesome5}
\usepackage{titlesec}
\usepackage{academicons}

% % --- Font: Source Sans Pro via Type1 (works with pdfLaTeX)
% \usepackage[default,scale=1.0]{sourcesanspro} % body + sans by default
% \renewcommand{\familydefault}{\sfdefault}      % make sans the default family
% \usepackage[scaled]{helvet}
\usepackage[default,scale=0.95]{opensans}

% (optional) math that pairs reasonably with sans body
\usepackage{newtxsf} % sans math; or use newtxtext/newtxmath for serif math

\usepackage[
  backend=biber,
  style=ieee,      % or 'ieee' / 'apa' / 'numeric-comp' depending on preference
  sorting=ydnt,          % Year (descending), name, title
  maxbibnames=99,        % show all authors instead of "et al."
  maxcitenames=2,        % optional: limit names in in-text citations
  giveninits=true,      % show full first names (set true for initials)
  uniquename=false,      % don't force initials to distinguish same names
  doi=false,             % suppress DOI (optional)
  url=false              % suppress URL unless needed
  ]{biblatex}
  \addbibresource{references.bib}
  
\usepackage{xstring}

\DeclareNameFormat{author}{%
\edef\fullname{\namepartgiven\space\namepartfamily}%
\IfBeginWith{\fullname}{P}{%
\IfSubStr{\fullname}{Popov}{%
\mkbibbold{\namepartfamily\addcomma\space\namepartgiven}%
}{%
\namepartfamily\addcomma\space\namepartgiven%
}%
}{%
\namepartfamily\addcomma\space\namepartgiven%
}%
\ifnum\value{listcount}<\value{liststop}
\addcomma\space
\fi
}

\DeclareBibliographyDriver{article}{%
\printnames{author}%
\setunit{\addcomma\space}%
\printfield[title]{title}%
\setunit{\addcomma\space}%
\printfield{journaltitle}%
\setunit{\addcomma\space}%
\printfield{year}%
\finentry
}

\DeclareBibliographyDriver{inproceedings}{%
\printnames{author}%
\setunit{\addcomma\space}%
\printfield[title]{title}%
\setunit{\addcomma\space}%
\printfield{booktitle}%
\setunit{\addcomma\space}%
\printfield{year}%
\finentry
}

\titleformat{\section}{\Large\bfseries}{\thesection}{1em}{}[\titlerule]
\titlespacing{\section}{0pt}{12pt}{6pt}

\usepackage{comicneue}

\begin{document}

\begin{center}
  % {\Huge\textbf{Pavel Popov}}\\[0.5em]
  % {\comicneue\Huge\textbf{Pavel Popov}}\\[0.5em]
  {\fontspec{JetBrains Mono}\Huge Pavel Popov}\\[0.5em]
  {\fontspec{JetBrains Mono}\footnotesize +1 (470) 920-1876 \quad | \quad Center street NW, Atlanta, GA 30318} \quad | \quad {paavali.popov@gmail.com}\\
  {\small \href{https://github.com/paavalipopov}{Github: paavalipopov} \quad | \quad
  % \href{https://www.linkedin.com/in/pavel-popov-83b69a123/}{\faLinkedin\ Pavel Popov} \quad
  \href{https://www.linkedin.com/in/pavel-popov-83b69a123/}{LinkedIn: pavel-popov-83b69a123} \quad | \quad
  \href{https://scholar.google.com/citations?user=4alUD5UAAAAJ}{Google scholar: 4alUD5UAAAAJ}}
\end{center}

\section*{Summary}
{\small \noindent PhD candidate in Computer Science with domain knowledge in physics and neuroscience, working with large multimodal health-related datasets in collaboration with cross-functional teams.
Strong programming and data analysis skills, proficiency with using various shell scripting, Python, PyTorch, MATLAB, MySQL and mongoDB, HPC clusters and distributed data parallel training, proven track record of interdisciplinary collaboration and research.

\noindent Seeking applied ML engineer/researcher/intern roles. 
Over 10 years of experience in independent research design, execution, analysis, and writing across different fields, with deliverables including 9 peer-reviewed publications, and oral presentation experience at US and international conferences.}

\section*{Skills}
\begin{itemize}
[leftmargin=*,noitemsep] 
\item \textbf{Programming \& Tools:} Python, PyTorch, Scikit-learn, Pandas, MATLAB, Git, Docker, HPC, SLURM, Spark, SQL, mongoDB, Figma.
\item \textbf{Machine Learning:} Deep learning architecture design, time series analysis, causal inference, explainable AI, distributed data parallel training.
\item \textbf{Data mining:} Frequent pattern analysis, feature engineering, data visualization.
% \item \textbf{Domain Expertise:} Neuroimaging, solid state physics, signal processing.
\end{itemize}

\section*{Education}
\noindent\textbf{Georgia State University} \hfill June 2027\\
PhD in Computer Science \hfill
% Machine Learning \& Neuroimaging, 
GPA 3.89\\[0.3em]
\textbf{Moscow Institute of Physics and Technology} \hfill June 2020\\
MS in Applied Mathematics and Physics \hfill 
% Spintronics, 
GPA 3.92\\[0.3em]
\textbf{Moscow Institute of Physics and Technology} \hfill June 2018\\
BS in Applied Mathematics and Physics \hfill 
% Solid State Physics, 
GPA 3.81

\section*{Experience}
\noindent\textbf{TReNDS Center} \hfill January 2022 --- Now\\
Graduate Research Assistant \hfill \textit{Atlanta, USA}\\
[-2em]
\begin{itemize}[leftmargin=*,noitemsep]
    \item Developed a robust DL classifier for medical diagnostics of functional MRI scans with SOTA performance and 10x speedup to existing benchmarks.
    \item Designed an interpretable AI model to simulate multivariate time-series data, enabling the identification of novel biomarkers and data-driven insights.
    \item Built and optimized frameworks for training and deployment of DL models for neuroimaging analysis in distributed systems.
    \item Coordinated work of cross-functional teams to develop ML tools useful in clinical practice.
    \item Presented and translated technical findings for general audiences, organized and volunteered in social events.
\end{itemize}

\noindent\textbf{Institute of Radioengineering and Electronics} \hfill May 2017 --- September 2020\\
Research Assistant \hfill \textit{Moscow, Russia}\\
[-2em]
\begin{itemize}[leftmargin=*,noitemsep]
    \item Modeled complex signal propagation in novel hardware architectures based on wave logic.
    \item Developed mathematical models for control of high-frequency signal dynamics, providing theoretical validation for prototype hardware controllers.
    \item Ran large-scale numerical simulations to validate theoretical models.
\end{itemize}

\newpage

\section*{Open Source Projects}
\noindent\textbf{\href{https://pypi.org/project/ml4fmri/}{ml4fmri}: Python toolkit for benchmarking DL models for fMRI analysis}\\
[-2em]
\begin{itemize}[leftmargin=*,noitemsep]
    \item Unified 13 DL models for multivatiate time series analysis under a standardized training API.
    \item Automated an all-model benchmarking pipeline for any input time series for reproducible research.
\end{itemize}

\noindent\textbf{\href{https://pypi.org/project/brainbow/}{brainbow} — Neuroimaging visualization utility}\\
[-2em]
\begin{itemize}[leftmargin=*,noitemsep]
    \item Created a command-line tool for visualizing obscure neuroimaging data formats.
\end{itemize}


\section*{Conferences \& Workshops}
\newcommand{\confattend}[4][1.2em]{%
\noindent\begin{tabular*}{\linewidth}{@{\extracolsep{\fill}} l r}
  \hspace*{#1}\textbf{#2} & #4 \\
  \hspace*{#1}\emph{#3} &
\end{tabular*}\par\vspace{0.25em}
}

\confattend{GSU Research Center Brain Health Synergy Symposium}{Atlanta, GA, USA}{Dec 2025}
\confattend{IEEE Workshop on Machine Learning for Signal Processing (MLSP 2025)}{Istanbul, Turkey}{Sep 2025}
\confattend{Organization for Human Brain Mapping (OHBM 2025)}{Brisbane, Australia}{Jun 2025}
\confattend{IEEE Brain Discovery and Neurotechnology Workshop}{Chicago, IL, USA}{Oct 2024}
\confattend{Organization for Human Brain Mapping (OHBM 2023)}{Montreal, Canada}{Jul 2023}
\confattend{Machine Learning for Science \& Engineering Workshop}{Seattle, WA, USA}{Jun 2023}
\confattend{Euro-Asian Symposium «Trends in Magnetism» (EASTMAG 2019)}{Ekaterinburg, Russia}{Sep 2019}
\confattend{International Congress on Ultrasonics (ICU 2019)}{Brugge, Belgium}{Sep 2019}




\section*{Publications}
\nocite{meanMLP}
\nocite{PhysRevAppl2020}
\nocite{JMMM2019}
\nocite{RiEl2018}
% \nocite{Safin2020}
% \nocite{Safin2020a}
\nocite{zafar}
\nocite{decifra_mlsp}
\nocite{vanduyn}


\printbibliography[
heading=none, 
sorting=ydnt
]

\end{document}
